\begin{frame}{데이터 고르기 실전 가이드 — 기준 (6)}
  \begin{itemize}
    \item Ch8.1의 5개 기준(행의 물리적 단위, 레이아웃 함정, 시계열,
      클래스 비율, 크기)은 \textbf{그대로} 적용
    \item Block B에서만 추가 — \kb{(6) ``원시 특징''이 이미 정형으로
      추출되어 있는가?} (정형/비정형 판정의 실전판)
  \end{itemize}
  \vskip0.4em
  \small
  \setlength{\tabcolsep}{4pt}
  \begin{tabular}{@{}p{2.1cm}p{5.9cm}p{5.8cm}@{}}
    \toprule
    경우 & 예 & 판정 \\
    \midrule
    (1) 수치·범주 열로만 구성된 표 & 매출표, 센서 로그(UCI 대부분) &
      \textbf{정형} — GBDT와 비교해야 한다 \\
    (2) 픽셀·음파·토큰인 원시 배열 & 32$\times$32$\times$3 이미지,
      오디오, 문장 & \textbf{비정형} — ``손으로'' 특징 만들기 어려우니
      \textbf{신경망이 특징 추출기} \\
    (3) 애매한 경우 & 이미 워드임베딩으로 벡터화된 텍스트 &
      ``원시 토큰으로 직접 모델링''을 선택해 Ch13 도구 필요성을
      스스로 만드는 것도 합법 — \textbf{이유를 M1에 적어야 한다} \\
    \bottomrule
  \end{tabular}
\end{frame}
